\chapter{Einleitung, Motivation und Zielsetzung}
\label{chap:einleitung}

Die Digitalisierung hat in den vergangenen Jahrzehnten nahezu alle Lebensbereiche verändert, von der
Arbeitswelt über die Kommunikation bis hin zu der Art, wie Wissen erworben und weitergegeben wird. Mit
dem Aufkommen generativer \gls{ki} hat sich diese Entwicklung in den letzten Jahren noch einmal
deutlich beschleunigt: Systeme, die auf natürlichsprachliche Eingaben reagieren und komplexe Texte,
Programmcode oder Argumentationen erzeugen können, stehen heute praktisch jeder Person mit
Internetzugang zur Verfügung, ohne dass technisches Vorwissen erforderlich wäre.

Auch die Hochschulbildung bleibt von dieser Entwicklung nicht unberührt, und kaum ein Bereich hat sich
dabei so schnell verändert wie sie. Innerhalb weniger Jahre haben sich Systeme wie ChatGPT, Microsoft
Copilot oder Gemini von einer Randerscheinung zu einem alltäglichen Begleiter für Millionen
Studierender weltweit entwickelt \cite{kasneciChatGPTGoodOpportunities2023, labadzeRoleAIChatbots2023}.
Auch an deutschen Hochschulen bestätigen aktuelle Studien eine bereits sehr hohe und weiter wachsende
Nutzung generativer \acrshort{ki} durch Studierende. Eine bundesweite Längsschnittbefragung von 4910
Studierenden an 395 Hochschulen weist einen Anstieg von 63,2 Prozent im Jahr 2023 auf 91,6 Prozent im
Jahr 2025 aus \cite{vongarrelKuenstlicheIntelligenzStudium2025,
marczukKuenstlicheIntelligenzStudienalltag2025}. Damit rückt eine grundlegende Frage in den Vordergrund der hochschuldidaktischen
Diskussion: Verändert generative \acrshort{ki}, wie Studierende lernen, Probleme durchdringen und
Entscheidungen treffen, und wenn ja, in welche Richtung \cite{limGenerativeAIFuture2023}?

Hochschulen begegnen dieser Entwicklung bislang vor allem mit Regularien und Richtlinien zum
verantwortungsvollen Umgang mit generativer \acrshort{ki}
\cite{chanComprehensiveAIPolicy2023, hallGenerativeAIReweaving2024}. Wie ein \acrshort{ki}-System
selbst so gestaltet werden kann, dass es eigenständiges Denken fördert statt zu ersetzen, wird
demgegenüber, insbesondere für offene Aufgaben ohne eindeutig richtige Lösung, noch seltener
systematisch untersucht (vgl. Abschnitt~\ref{sec:standtechnik}). Diese gestalterische Frage lässt sich kaum
abstrakt beantworten, sondern nur an konkreten Aufgabenstellungen untersuchen, die tatsächlich eine
eigene Begründung und Abwägung verlangen und nicht mit einer einzigen richtigen Antwort auskommen.

Diese Arbeit setzt genau an dieser Frage an, und zwar im Software Engineering, einem Fachgebiet, in
dem Studierende bereits im Studienalltag intensiv mit \acrshort{ki}-gestützten Werkzeugen arbeiten und
in dem Prozessentscheidungen, wie im weiteren Verlauf dieses Kapitels ausgeführt wird, eine besonders
geeignete Aufgabenklasse darstellen, um Reflexionstiefe überhaupt sichtbar zu machen.

Dieses Kapitel legt zunächst die Motivation für die vorliegende Arbeit dar
(Abschnitt~\ref{sec:motivation}), leitet daraus die konkrete Problemstellung ab
(Abschnitt~\ref{sec:problemstellung}) und formuliert abschließend die Zielsetzung sowie die daraus
abgeleiteten Forschungsfragen (Abschnitt~\ref{sec:zielsetzung}).

Die weitere Arbeit ist wie folgt aufgebaut: Kapitel~\ref{chap:theogrundlage} führt die technischen
und didaktischen Grundlagen sowie den Stand der Technik ein und arbeitet die Forschungslücke heraus,
an der diese Arbeit ansetzt. Kapitel~\ref{cha:methodik} beschreibt das Untersuchungsdesign, die Datenerhebung
und die Datenanalyse. Es steht bewusst vor der Konzeptionierung, weil es den Auswertungsrahmen
festlegt, auf den sich die anschließenden Gestaltungsentscheidungen beziehen.
Kapitel~\ref{cha:concept} dokumentiert die Konzeption und technische Umsetzung des Reflexionsbots.
Kapitel~\ref{cha:ergebnisse} berichtet die Ergebnisse rein darstellend, bevor
Kapitel~\ref{cha: diskussion} sie interpretiert, die Forschungsfragen beantwortet, die Befunde in den
Forschungsstand einordnet und Limitationen sowie einen Ausblick darlegt.
Kapitel~\ref{chap:fazit} fasst die zentralen Erkenntnisse zusammen.

\section{Motivation}
\label{sec:motivation}

Studierende nutzen generative \acrshort{ki}-Systeme längst nicht mehr nur zur Informationssuche,
sondern aktiv bei der Bearbeitung komplexer Aufgabenstellungen, auch im Rahmen von Lehrveranstaltungen
\cite{nikolicChatGPTCopilotGemini2024, jwairGenerativeArtificialIntelligence2025}. Diese Praxis wird
kontrovers beurteilt. Generative \acrshort{ki} kann Lernprozesse durch personalisierte Unterstützung
bereichern \cite{baidoo-anuEducationEraGenerative2023}, wird jedoch auch genutzt, um Aufgaben schnell
zu erledigen, ohne sich inhaltlich mit den Lerninhalten auseinanderzusetzen
\cite{francisGenerativeAIHigher2025, cottonChattingCheatingEnsuring2024}. Eine Meta-Analyse zum
kognitiven Einfluss generativer \acrshort{ki}-Werkzeuge legt nahe, dass die Wirkung weniger davon abhängt,
ob \acrshort{ki} überhaupt eingesetzt wird, als davon, wie ein System gestaltet und in die
Lehrveranstaltung eingebettet ist \cite{GenerativeAITools}. Hochschulen richten ihr Augenmerk deshalb
zunehmend darauf, wie ein System gestaltet sein muss, damit es eigenständiges Denken fördert statt zu
ersetzen \cite{chanComprehensiveAIPolicy2023}.

Im Software Engineering stellt sich diese Frage mit besonderer Dringlichkeit, da Studierende hier
bereits im Studienalltag intensiv mit \acrshort{ki}-gestützten Werkzeugen arbeiten
\cite{sengulSoftwareEngineeringEducation2024}. Erste Erfahrungsberichte zu \acrshort{ki}-Tutoren in
der Software-Engineering-Lehre lassen offen, ob solche Systeme zu einem tieferen Verständnis führen
oder vor allem den Bearbeitungsaufwand reduzieren
\cite{frankfordAITutoringSoftwareEngineering2024, happeExperienceReportPedagogically2025}. Diesem
Risiko soll unter anderem durch sokratisches Nachfragen anstelle direkter Lösungsvorschläge begegnet
werden. Bislang wurde dieser Ansatz jedoch überwiegend für eng abgegrenzte Aufgaben, etwa einzelne
Programmierfehler, untersucht \cite{karguptaInstructNotAssist2024}.

Besonders deutlich wird dieses Spannungsfeld im
Software-Prozessmanagement, wo sich Prozessentscheidungen selten mit einer einzigen richtigen
Antwort lösen lassen, sondern eine Abwägung konkurrierender Ziele wie Geschwindigkeit,
Dokumentationsaufwand und Compliance erfordern. Reflexion über die eigenen Entscheidungen gilt dabei,
wie Befunde aus der Entscheidungsdidaktik in anderen Fachkontexten zeigen, als zentral für den Aufbau
von Entscheidungskompetenz \cite{greschEnhancingDecisionMakingSTSE2017}. Ein
generisches \acrshort{ki}-System, das auf Wunsch eine vollständige, plausibel klingende Lösung liefert,
kann diesen Reflexionsprozess umgehen, ohne dass dies für Studierende oder Lehrende sichtbar wird.
Motiviert wird diese Arbeit deshalb von der Frage, ob ein \acrshort{ki}-System, das gezielt für
Reflexion statt für schnelle Antworten gestaltet ist, Studierende anders bei
Prozessentscheidungen unterstützt als ein generisches System, mit dem sie ohnehin bereits im
Studienalltag arbeiten.

\section{Problemstellung}
\label{sec:problemstellung}

Für die didaktische Gestaltung generativer \acrshort{ki} ist entscheidend, wie ein System auf
Eingaben reagiert. Generische Systeme sind darauf ausgelegt, auf jede Eingabe eine möglichst
hilfreiche, direkte Antwort zu liefern. Für einfache Wissensfragen ist das unproblematisch, für
Aufgaben, die eine eigene Begründung und Abwägung verlangen, besteht jedoch die Gefahr, dass
Studierende die gelieferte Antwort ungeprüft übernehmen, ohne die zugrunde liegenden Kompromisse
selbst durchdacht zu haben \cite{cottonChattingCheatingEnsuring2024}.

Problematisch ist dabei weniger die Existenz einer Antwort an sich als vielmehr ihre Wirkung auf den
Bearbeitungsprozess: Eine vollständige, kohärent formulierte Lösung wirkt überzeugend, unabhängig
davon, ob sie durch eigenes Abwägen oder durch bloße Übernahme zustande gekommen ist. Für Lehrende ist
im fertigen Leistungsprodukt kaum erkennbar, welcher der beiden Wege tatsächlich beschritten wurde.
Diese Unsichtbarkeit macht es schwer, allein anhand der Ergebnisse zu beurteilen, ob ein
\acrshort{ki}-Einsatz die intendierten Lernprozesse tatsächlich unterstützt oder lediglich verdeckt,
dass sie ausgeblieben sind.

Alternative Ansätze, die Antworten bewusst zurückhalten und stattdessen durch Rückfragen zur
Auseinandersetzung anregen, existieren bereits, wurden jedoch bislang vor allem für eng abgegrenzte,
klar strukturierte Aufgaben wie die Suche nach einem einzelnen Programmierfehler untersucht
\cite{karguptaInstructNotAssist2024}. Ob ein solcher Ansatz auch in einem Themenfeld funktioniert, in
dem es keine eindeutig richtige Lösung gibt, sondern mehrere vertretbare Optionen mit unterschiedlichen
Konsequenzen abzuwägen sind, wie es für Prozessentscheidungen im Software-Prozessmanagement typisch
ist, ist bislang kaum empirisch untersucht.

Es fehlt insbesondere an einem direkten, praktischen Vergleich zwischen einem generischen
KI-System und einem dediziert für Reflexion gestalteten Gegenstück unter realen Bedingungen einer
Lehrveranstaltung, der beide Systeme derselben Aufgabenstellung gegenüberstellt und die tatsächlich
gezeigte Reflexionstiefe anhand der entstandenen Bearbeitungen und Chatverläufe nachvollzieht, statt
sich allein auf die subjektive Einschätzung der Studierenden zu verlassen. Diese Arbeit setzt an dieser
Lücke an.

\section{Zielsetzung und Forschungsfragen}
\label{sec:zielsetzung}
\label{sec:forschungsfrage}

Erkenntnisleitendes Ziel dieser Arbeit ist es zu klären, ob und in welcher Hinsicht sich die
didaktische Gestaltung eines KI-Systems darauf auswirkt, wie eigenständig und reflektiert Studierende
anspruchsvolle Prozessentscheidungen im Software Engineering bearbeiten. Nicht die Frage, ob ein
solches System technisch funktioniert, steht dabei im Vordergrund, sondern die Frage, ob ein gezielt
auf Reflexion statt auf schnelle Antworten ausgelegtes System Studierende anders unterstützt als ein
generisches KI-System, mit dem sie ohnehin bereits im Studienalltag arbeiten.

Daraus ergibt sich die zentrale Forschungsfrage dieser Arbeit:

\begin{quote}
Inwiefern unterscheidet sich ein didaktisch gestalteter Reflexionsbot von einem generischen
KI-System darin, wie eigenständig und reflektiert Studierende Prozessentscheidungen im
Software-Prozessmanagement bearbeiten?
\end{quote}

Sie wird in drei Teilfragen (TF1 bis TF3) ausdifferenziert, die jeweils an ein konkret erhobenes
Konstrukt gebunden sind:

\begin{itemize}
    \item[\textbf{TF1}] Wie unterscheidet sich zwischen beiden Systemen der Anteil an die KI
    ausgelagerter beziehungsweise ungeprüft übernommener Bearbeitungen? \emph{(Operationalisierung:
    Autorenschaftskodierung der Bearbeitungen, vgl. Abschnitt~\ref{sec:datenanalyse}.)}
    \item[\textbf{TF2}] Welche Reflexionstiefe und welches selbsteingeschätzte Prozessverständnis
    zeigen beziehungsweise berichten Studierende beim Reflexionsbot im Vergleich zum generischen
    System? \emph{(Operationalisierung: Reflexionsstufen nach Hatton und Smith sowie die
    Prä-Post-Items zu Beschreiben, Begründen und Abwägen, vgl.
    Abschnitte~\ref{sec:reflexionsstufenmodelle} und~\ref{sec:fragebogen-experiment}.)}
    \item[\textbf{TF3}] Wie bewerten Studierende ihre eigene Kompetenzentwicklung sowie die Stärken
    und Grenzen beider Ansätze für ihr Lernen? \emph{(Operationalisierung: Post-Survey-Items zu
    Kompetenzentwicklung, Erwartung und Nutzungspräferenz, vgl.
    Abschnitt~\ref{sec:fragebogen-experiment}.)}
\end{itemize}

Die Beantwortung dieser Fragen soll zeigen, ob und wie sich ein didaktisch gestalteter
Reflexionsbot in den erhobenen Daten von einem generischen System unterscheidet, und welche
Konsequenzen sich daraus für den gezielten KI-Einsatz in der Hochschullehre ergeben.

Als methodischer Weg zur Beantwortung dieser Fragen wurde ein Reflexionsbot konzipiert und
umgesetzt, der Studierende durch sokratisches Fragen statt direkter Lösungen begleitet. Um eine
Aufgabenstellung zu schaffen, an der sich Reflexionstiefe überhaupt sinnvoll unterscheiden lässt,
wurden zwei realitätsnahe Fallszenarien aus dem Software-Prozessmanagement konzipiert, die jeweils
mehrere Dilemma-Situationen mit widerstreitenden Zielen enthalten, etwa Patientensicherheit
gegenüber Entwicklungsgeschwindigkeit oder Datenschutzanforderungen gegenüber Projektfortschritt.
Details zur Gestaltung dieser Szenarien und des Reflexionsbots folgen in Kapitel~\ref{cha:concept}.

Der Reflexionsbot wurde im Rahmen eines Mini-Experiments in der Lehrveranstaltung \glqq Grundlagen
des Software Engineering 1\grqq{} im ersten Fachsemester des Bachelorstudiengangs Software
Engineering (B.Sc.) an der Hochschule Heilbronn eingesetzt und einer Vergleichsgruppe
gegenübergestellt, die dieselben Szenarien mit einem generischen KI-System bearbeitete. Die Auswertung erfolgt als
Mixed-Methods-Untersuchung: Quantitativ werden Prä- und Post-Survey-Daten zur wahrgenommenen Wirkung
beider Systeme ausgewertet, qualitativ werden die tatsächlichen schriftlichen Bearbeitungen und
Chatverläufe mittels strukturierender Inhaltsanalyse auf ihre Reflexionstiefe hin untersucht. Erst
die Zusammenführung beider Perspektiven erlaubt es, wahrgenommene und tatsächlich gezeigte Reflexion
zu unterscheiden, statt sich allein auf die Selbsteinschätzung der Studierenden zu verlassen.

\section{Abgrenzung des Untersuchungsrahmens}
\label{sec:abgrenzung}

Der Untersuchungsrahmen dieser Arbeit ist bewusst eng gefasst, um den Vergleich beider Systeme unter
kontrollierten und im Rahmen einer Bachelorarbeit umsetzbaren Bedingungen durchführen zu können.

Gegenstand ist ein einzelner Anwendungsfall, nämlich Prozessentscheidungen im
Software-Prozessmanagement, bearbeitet anhand von zwei eigens konstruierten Fallszenarien in einer
Übungseinheit einer einzelnen Lehrveranstaltung. Diese Fokussierung erlaubt es, beide Gruppen mit identischen Aufgabenstellungen zu
konfrontieren und die entstandenen Bearbeitungen vollständig auszuwerten, statt Vergleichbarkeit
zugunsten einer breiteren Abdeckung aufzugeben.

Nicht Gegenstand der Arbeit sind daher die Übertragbarkeit auf andere Fachgebiete, Kohorten oder
Hochschulen, ein längerfristiger Lernzuwachs über die Übungseinheit hinaus sowie ein kontrollierter
Vergleich verschiedener Sprachmodelle. Der Reflexionsbot wird mit einem einzigen Modell umgesetzt,
während die Vergleichsgruppe ein frei gewähltes, in seiner Modellstufe nicht kontrolliertes
Consumer-System nutzt (vgl. Abschnitt~\ref{sec:fragebogen-experiment}). Die Auswertung erfolgt
angesichts der kleinen Stichprobe ausschließlich deskriptiv und explorativ. Auf inferenzstatistische
Tests und Verallgemeinerungsansprüche wird bewusst verzichtet. Welche Konsequenzen sich daraus für die
Reichweite der Ergebnisse ergeben, wird in Abschnitt~\ref{sec:limitationen} ausführlich dargelegt,
mögliche Erweiterungen benennt der Ausblick in Abschnitt~\ref{sec:ausblick}.

\section{Beitrag der Arbeit}
\label{sec:beitrag}

Aus dieser Zielsetzung ergeben sich vier Beiträge:

\begin{description}
    \item[Ein dokumentierter, reproduzierbarer Reflexionsbot.] Der entwickelte Bot ist nicht nur
    konzeptionell beschrieben, sondern mit eingesetztem Sprachmodell, Architektur und der
    Systemnachricht im vollständigen Wortlaut dokumentiert (Kapitel~\ref{cha:concept},
    \hyperref[app:systemprompt]{Anhang~A.7}), sodass sich der didaktische Eingriff nachvollziehen und replizieren lässt.
    \item[Ein zweidimensionales Kategoriensystem.] Die qualitative Auswertung erfasst neben der
    Reflexionsstufe nach Hatton und Smith die \emph{Autorenschaft} der eingereichten Texte als eigene
    Dimension mit prüfbaren Indikatoren (\hyperref[app:kategoriensystem]{Anhang~A.1}). Damit wird verhindert, dass die Analyse die
    Reflexionsfähigkeit des KI-Systems statt der Studierenden misst, ein Problem, das bei offenen
    Aufgaben mit KI-Zugang unmittelbar auftritt.
    \item[Empirische Evidenz für offene Prozessentscheidungen.] Der Vergleich zwischen
    spezialisiertem und generischem System wird auf Aufgaben ohne eindeutig richtige Lösung
    übertragen, ein Kontext, für den ein solcher Vergleich bislang fehlt
    (Abschnitt~\ref{sec:standtechnik}).
    \item[Ein Beleg für den Wert der Triangulation.] Der fallweise Abgleich von Selbstauskunft und
    dokumentiertem Bearbeitungsverhalten zeigt, dass beide auseinanderfallen können und eine
    Befragung allein die tatsächliche Auseinandersetzung nicht zuverlässig abbildet
    (Abschnitt~\ref{sec:selbstauskunft-verhalten}).
\end{description}